\Needspace{10\baselineskip}
\begin{origexamplebox}{例题 21.1【卡诺热机】}
利用熵增原理证明 Carnot 定理：任意工作物质作 Carnot 循环的效率都只与高、低温热库的温度有关，
与工质的数量、物态方程均无关，效率均为
\[
  \eta=1-\frac{T_2}{T_1},
\]
其中 $T_1,T_2$ 分别为高、低温热库的温度。
\end{origexamplebox}

\begin{xesolution}
设热机在一个可逆循环中从温度为 $T_1$ 的高温热库吸收热量 $Q_1$，向温度为 $T_2$ 的低温热库放出
热量 $Q_2$。工质经历完整循环后回到初态，作为状态函数的熵没有净改变。

高温热库失去热量，熵变为 $-Q_1/T_1$；低温热库得到热量，熵变为 $Q_2/T_2$。Carnot 循环是可逆
循环，把工质、两热库和只储存机械能的理想功库一并作为孤立整体；功库熵不变，所以总熵变为零，因此
\[
  -\frac{Q_1}{T_1}+\frac{Q_2}{T_2}=0,
  \qquad
  \frac{Q_2}{Q_1}=\frac{T_2}{T_1}.
\]

\carnotProofDiagram
\figcap{卡诺热机的能量流与熵平衡。$Q_1,Q_2,W$ 均表示图中箭头方向上的正量。}\label{fig:carnot-proof}
图\ref{fig:carnot-proof}中热量的大小不同，但可逆情况下从高温热库取走的熵等于送给低温热库的熵。

热机一个循环的对外净功为 $W=Q_1-Q_2$，故
\[
  \boxed{\eta=\frac{W}{Q_1}
  =1-\frac{Q_2}{Q_1}
  =1-\frac{T_2}{T_1}}.
\]
推导中没有使用工质的物态方程、热容量或物质的量，只使用了“循环后工质熵不变”和“整个过程可逆”。
所以任意工质作 Carnot 循环都具有同一效率。若过程不可逆，总熵会增加，向低温热库排出的热量更多，
实际效率必低于 Carnot 效率。
\end{xesolution}

\Needspace{9\baselineskip}
\begin{origexamplebox}{例题 21.2【熵增原理】}
利用熵增原理，重新对热力学第二定律的 Kelvin 和 Clausius 陈述给出定量说明。
\end{origexamplebox}

\begin{xesolution}
\textbf{Kelvin 表述。} 假设一台循环装置从温度为 $T$ 的单一热库吸收热量 $Q$，并把它全部变为功。
装置完成循环后熵变为零；热库失去热量，熵变为
\[
  \Delta S_{\rm reservoir}=-\frac{Q}{T}<0.
\]
还需把接收机械功的理想功库纳入整体，例如无摩擦地被提高的重物；它只增加机械能，熵不变。
于是孤立的“装置、热库加功库”总熵减少，违背熵增原理。因此不可能在没有其他变化的条件下，从单一热库
吸热并把热量全部转化为功。

\textbf{Clausius 表述。} 假设一台循环装置不消耗功，却把热量 $Q$ 从温度 $T_c$ 的低温热库搬到温度
$T_h$ 的高温热库，其中 $T_h>T_c$。两热库的总熵变为
\[
  \Delta S_{\rm tot}
  =\frac{Q}{T_h}-\frac{Q}{T_c}<0.
\]
这同样违背熵增原理。因此把热量从低温处送往高温处，必须伴随外界做功或其他补偿变化。

\secondLawEntropyDiagram
\figcap{两种被禁止的唯一效果及其负熵变。比较时都须把装置和所有环境变化计入。}\label{fig:second-law-comparison}
图\ref{fig:second-law-comparison}左侧的热机必须已完成循环；单次等温膨胀并不属于这个禁止条件。

两个陈述表面上分别谈“热变功”和“逆向传热”，本质上都排除了使孤立系统总熵减少的循环装置。
\end{xesolution}

\Needspace{10\baselineskip}
\begin{origexamplebox}{例题 21.3【统统证出来】}
请利用热力学第二定律辨析下列说法，指出所需条件，并给出正确结论：
\begin{enumerate}[label=\alph*.]
  \item 理想气体的等温膨胀过程是不可能自然发生的；
  \item 对本讲的定量理想气体，可逆绝热线与等温线不能相交两次；
  \item 对固定组成、稳定单相的简单可压缩系统，两条不同可逆绝热线不能相交；
  \item 有限温差热传递过程不可逆；
  \item 快速膨胀过程不可逆。
\end{enumerate}
\end{origexamplebox}

\begin{xesolution}
这组判断的共同方法是：若假定某过程可以在没有其他影响的情况下逆转，就尝试拼出第二类永动机，或
直接算出总熵减少。

\textbf{a. 等温膨胀。} 这一说法本身不成立，不能给“自然发生”附加原本没有的限制来证明它。
理想气体在恒温环境中推动活塞缓慢膨胀可以发生，单个过程并不违反第二定律。此时 $\Delta U=0$，对外功
\[
  W=nRT\ln\frac{V_2}{V_1}>0
\]
必须完全来自热库供热。工质的体积已经改变，这不是“装置循环复原后只留下热变功”这一被禁止的效果。
它不能在既绝热又输出功的条件下仍维持等温，这首先可由第一定律判断。另一个更直接的反例是撤去
隔板，让理想气体向真空自由膨胀，则 $Q=W=\Delta U=0$，温度确实不变，而且过程会自发发生；这正是
不可逆等温自由膨胀，不能与可逆等温膨胀混为一谈。

\irreversibleProcessesDiagram
\figcap{左侧分别画出撤去隔板前、后的气体状态；右侧为两恒温热库间的有限温差传热。}\label{fig:irreversibility}
图\ref{fig:irreversibility}中左图初态有真空区，末态虚线只标隔板原位置，不再是阻挡气体的实体壁；
右图的传热将在第d问中计算总熵变。

\textbf{b. 绝热线与等温线的交点。} 对本讲的理想气体，固定温度下
$S(T,V)=\text{常量}+nR\ln V$ 随 $V$ 严格增加，而同一可逆绝热线有固定熵，故只能选出一个 $V$。
也可代入 $P=nRT/V$ 与 $PV^\gamma=K$，得 $nRTV^{\gamma-1}=K$，在 $\gamma>1,V>0$ 时解唯一。
若强行假设两次相交并能围成非零功的循环，就可用一段等温线与一段绝热线构造单热库循环，违反第二定律。
对于一般稳定单相系统，还需可逆绝热线上的温度随体积单调，例如通常热膨胀系数为正的介质。
仅有正热容不能保证这个单调性；存在反常热膨胀时，不能把“两次相交不可能”无条件推广。

\textbf{c. 两条绝热线不能相交。} 可逆绝热线就是等熵线。若两条绝热线在某平衡状态相交，它们在该点
具有相同的熵；对于给定组成、单相稳定且物态方程光滑的简单系统，$c_v>0$ 使
$(\partial S/\partial T)_V=nc_v/T>0$，故固定熵和体积可唯一确定温度。通过该点的可逆绝热路径由微分方程唯一确定，因此
两条曲线实际是同一条绝热线，而不是两条不同曲线。若把它们看作不同曲线，还可以在两曲线间补上一小段
等温线，构造出上一问所排除的单热库循环。

\textbf{d. 有限温差传热不可逆。} 先设两个恒温热库，热量 $Q$ 从 $T_h$ 传到 $T_c$，$T_h>T_c$。总熵变
\[
  \Delta S_{\rm tot}
  =-\frac{Q}{T_h}+\frac{Q}{T_c}
  =Q\left(\frac1{T_c}-\frac1{T_h}\right)>0.
\]
若要完全逆转，就必须使孤立系统熵减少，故不能自发实现。只有在传热温差趋于零的可逆极限中，熵产生
才趋于零。
若两个有限物体的温度会变化，则应对微小传热量积分：
$\Delta S_{\rm tot}=\int(1/T_c-1/T_h)\,\delta Q>0$；不能把初始温度直接代入有限热量公式。

\textbf{e. 快速膨胀。} “快”本身不是可直接代入熵公式的物理条件。通常快速膨胀来不及维持平衡，
伴随有限压差、黏滞、冲击或最终弛豫时，会产生熵，因而不可逆；若只给一个短时间而没有机制，证明并不充分。
以绝热容器中撤去隔板、气体向真空自由膨胀为明确模型：气体由 $V_i$ 扩展到 $V_f$，温度不变，但
\[
  \Delta S=nR\ln\frac{V_f}{V_i}>0.
\]
若想让气体与环境都无变化地回到初态，就必须使总熵减少，所以过程不可逆。

\end{xesolution}

\Needspace{9\baselineskip}
\begin{origexamplebox}{例题 21.4【热传递】}
三个热容量均为 $c$ 的物体，已知 $T_1=T_2=\SI{300}{K}$，$T_3=\SI{100}{K}$。若外界不做功也不供热，
利用热机工作于三物体之间，使三物中某一物体温度升高，求它所能达到的最高温度；此时另外两物体的
温度各为多少？
\end{origexamplebox}

\begin{xesolution}
先证明任何允许过程的温度上界，再证明这个上界可由可逆装置趋近，不能只用两条等式猜测最优状态。
取最终最高温度为 $X$，另外两体温度为 $Y,Z$，均以开尔文计。热机完成循环，不储存净能量，
则第一定律和熵增原理分别给出
\[
 X+Y+Z=700,\qquad
 \Delta S=c\ln\frac{XYZ}{300^2\times100}\ge0.
\]
所以 $XYZ\ge9.0\times10^6$。又由 $(Y-Z)^2\ge0$，
\[
 YZ\le\left(\frac{Y+Z}{2}\right)^2=\left(\frac{700-X}{2}\right)^2,
 \quad X(700-X)^2\ge36.0\times10^6.
\]
因为 $X$ 是最高温度，$X\ge700/3$。函数 $f(X)=X(700-X)^2$ 在该区间内有
$f'(X)=(700-X)(700-3X)\le0$；而 $f(400)=36.0\times10^6$，故必有 $X\le400$。
等号同时要求总熵不变及 $Y=Z$。这才说明可逆性和另外两体等温是上界成立的必要条件。

设目标物体的最高温度为 $T_{\max}$。三个物体的热容量均为常量 $c$，系统与外界既不交换热量也不交换
功，因此总能量守恒：
\[
  c(T_{\max}+2T)=c(300+300+100),
  \qquad T_{\max}+2T=700.
\]
可逆极限下总熵不变：
\[
  c\ln\frac{T_{\max}}{300}
  +c\ln\frac{T}{300}
  +c\ln\frac{T}{100}=0.
\]
于是
\[
  T_{\max}T^2=300\times300\times100.
\]

\threeBodiesDiagram
\figcap{两台可逆装置在三物体之间协同工作：热机的功全部驱动热泵，不从外界取得净功。}\label{fig:three-body-engine}

联立 $T_{\max}=700-2T$，得到
\[
  (700-2T)T^2=9.0\times10^6.
\]
物理解为 $T=\SI{150}{K}$，从而
\[
  \boxed{T_{\max}=\SI{400}{K}},
  \qquad
  \boxed{T_{\rm other}=\SI{150}{K},\ \SI{150}{K}}.
\]
方程还有 $T=\SI{300}{K},T_{\max}=\SI{100}{K}$ 这一根，它只是把初态中三物体的编号重新排列，并非
所求的最高温度状态。
最后验证可实现性。图\ref{fig:three-body-engine}中，令供热体 $Y$ 与低温体 $Z$ 之间的可逆热机输出功，
恰好驱动从 $Z$ 向目标体 $X$ 输热的可逆热泵。每小步都令两机完成循环，故没有外界净功输入。
热机从 $Y$ 取出 $-c\dd Y>0$，输出功
$\delta W=(-c\dd Y)(1-Z/Y)$；热泵向 $X$ 送入 $c\dd X$，需要功
$\delta W=c\dd X(1-Z/X)$。二者相等，等价于
\[
 \frac{\dd X}{X}+\frac{\dd Y}{Y}+\frac{\dd Z}{Z}=0,
 \qquad \dd X+\dd Y+\dd Z=0.
\]
沿这一可逆路径，$XYZ=9.0\times10^6$ 且 $X+Y+Z=700$ 始终成立。对 $300\le X<400$，
另外两个正根是
\[
 Y,Z=\frac{700-X\pm\sqrt{(700-X)^2-36.0\times10^6/X}}{2},
\]
并满足 $X\ge Y>Z$，热机与热泵均可按图运行。随着 $X\to400$，$Y,Z\to150$，
两者的温差趋于零，达到可逆极限；实际有限速率装置只能逼近这一最高温度。
所有对数都作用于温度比，温度积的两边单位都是 $\mathrm K^3$。
\end{xesolution}

\Needspace{10\baselineskip}
\begin{origexamplebox}{例题 21.5【液体表面能】}
小明学习了 Carnot 热机后，想把液体表面作为工作物质。液体表面能够对外界物体施加表面张力，因此
表面具有能量，并且可以得到
\[
  U=\left(\sigma-T\frac{\dd\sigma}{\dd T}\right)A,
\]
其中 $\sigma$ 为液体的表面张力系数，$A$ 为液体表面积。请研究利用液体表面作为工作物质的热机，并
求出其效率。
\end{origexamplebox}

\begin{xesolution}
先约定 $A$ 是所有被计入的气液界面的总面积；若是肥皂膜，须把两面都计入，不能再额外乘2。
本题 $U$ 指随界面增加而产生的表面过剩内能，不包括整块液体的体相内能。
可逆扩展表面时外界对表面做元功 $\delta W^{\rm in}=\sigma\dd A$，因而
\[
 \dd U=T\dd S+\sigma\dd A,\qquad
 \delta W^{\rm out}=-\sigma\dd A.
\]
定义亥姆霍兹（Helmholtz）自由能 $F=U-TS$，有 $\dd F=-S\dd T+\sigma\dd A$。
在忽略曲率效应、$\sigma$ 仅随温度变化时，表面过剩自由能为
\[
  F=\sigma(T)A.
\]
又因为 $F=U-TS$，把题给的 $U$ 代入可得表面熵
\[
  S=\frac{U-F}{T}
  =-A\frac{\dd\sigma}{\dd T}.
\]
普通液体的表面张力通常随温度升高而减小，故 $\dd\sigma/\dd T<0$；在给定温度下扩大表面会增加熵，
并需要从热库吸热。

让液体表面工作在温度分别为 $T_h,T_c$ 的两个热库之间，组成两段等温过程和两段等熵过程。在高温
$T_h$ 下，表面积从 $A_1$ 可逆增加到 $A_2$，表面吸热
\[
  Q_h=T_h\Delta S,
  \qquad
  \Delta S=-\left.\frac{\dd\sigma}{\dd T}\right|_{T_h}(A_2-A_1)>0.
\]
然后沿等熵过程降到 $T_c$，在低温下使表面积减小并向低温热库放热。由于两条等温段的熵变量大小相等，
放出的热量为
\[
  Q_c=T_c\Delta S.
\]
所以一个循环的净功与效率为
\[
  W_{\rm out}=Q_h-Q_c=(T_h-T_c)\Delta S,
  \qquad
  \boxed{\eta=1-\frac{T_c}{T_h}}.
\]

\surfaceEngineDiagram
\figcap{表面工质的可逆循环。上下为等温段，两条弯曲边为等熵段；面积与温度都随等熵过程改变。}\label{fig:surface-cycle}

还需确定等熵线，不能直接把它画成等面积直线。对 $S=-A\sigma'(T)$ 微分得
\[
 \dd S=-\sigma'(T)\dd A-A\sigma''(T)\dd T=0,
 \qquad A\sigma'(T)=\text{常量}.
\]
一般等熵过程既改变面积也改变温度。图\ref{fig:surface-cycle}采用一个明确的非退化模型
$\sigma=\sigma_*-aT^2$，其中 $a>0$，温区内 $\sigma>0$。这时
$S=2aAT$、$U=(\sigma_*+aT^2)A$、$C_A=2aAT>0$，等熵线为 $AT=\text{常量}$。
高温等温段 $1\to2$ 的面积为 $A_1\to A_2$；冷却到低温的状态3满足
$A_3=A_2T_h/T_c$；低温压缩到状态4时 $A_4=A_1T_h/T_c$，再沿等熵线回到状态1。
由此 $Q_h=2aT_h^2(A_2-A_1)$、$Q_c=2aT_hT_c(A_2-A_1)$，再次给出相同效率。

如果只在很窄温区用线性表面张力近似，形式上 $\sigma'$ 为常量，等熵线会变成 $A=\text{常量}$，
循环呈矩形，形式功为
\[
  W_{\rm out}=(\sigma_c-\sigma_h)(A_2-A_1).
\]
但仅含表面项时这还导致 $C_A=-TA\sigma''=0$，是退化近似，不能据此描述真实液膜如何绝热升降温。
若计入体相自由能 $F_0(T)$，总熵变成 $S_0(T)-A\sigma'(T)$，等熵线必须重新求，通常不再竖直。
无论哪一可靠模型，只要可逆且工作在两个热库之间，效率仍为 $1-T_c/T_h$。
在 $\sigma$--$A$ 图中，本题是逆时针循环，因 $W_{\rm out}=-\oint\sigma\dd A$，其净功为正；
这与气体的 $W_{\rm out}=\oint P\dd V$ 符号相反，不能只凭“循环方向”判断正负。
\end{xesolution}
